\section{Chương~\ref{sec:dendrites}. Số đuôi gai}

Cấu tạo của các lớp và số đầu vào đã được đo bằng giải phẫu và ghi điện; ước tính~\eqref{eq:dendriteload} là vật lý đơn giản của tụ điện, còn lời giải thích về tính toán cho số đầu vào dựa trên lý thuyết và mô hình.

{\footnotesize
\setlength{\tabcolsep}{3pt}\renewcommand{\arraystretch}{1.15}
\begin{longtable}{>{\raggedright}p{0.5cm}>{\raggedright}p{4.0cm}>{\raggedright}p{5.6cm}>{\raggedright}p{2.3cm}>{\raggedright\arraybackslash}p{1.7cm}}
\toprule
TT & Khẳng định & Thí nghiệm và điều đã đo & Nguồn & Trạng thái \\
\midrule\endfirsthead
\toprule
TT & Khẳng định & Thí nghiệm và điều đã đo & Nguồn & Trạng thái \\
\midrule\endhead
1 & Điện dung riêng của màng $c_m\approx1$~\textmu F/cm$^2$ & Màng được lấy từ thân nơ-ron vào pipet, đặt một bậc điện áp và từ sự suy giảm của dòng điện dung tìm ra điện dung toàn phần, rồi chia cho diện tích: $0.9$~\textmu F/cm$^2$ ở ba lớp nơ-ron & \cite{Gentet2000} & đã đo \\
2 & Thời gian nạp $t\approx c_mA\Delta V/I$~\eqref{eq:dendriteload}: nơ-ron lớn cần hàng chục đầu vào đồng thời, nơ-ron nhỏ chỉ cần một khớp lớn & Ước tính theo định luật nạp tụ; các con số ($20$ và $300$~pF, $0.1$ và vài nA) là bậc độ lớn & suy ra trong chương & lý thuyết \\
3 & Một khớp khổng lồ cho dòng vài nanoampe và đưa ngay nơ-ron nhỏ tới ngưỡng & Ghi đồng thời đầu tận cùng và tế bào nhận trong lát não: dòng của một khớp $2$--$13$~nA, mỗi xung đầu vào gây một đáp ứng trên ngưỡng, độ trễ khoảng $0.4$~ms ở $36^\circ$C; tế bào nhận xuất \nt{GLYC} & \cite{Borst1995,BorstSoria2012} & đã đo \\
4 & Không đuôi gai: ở nơ-ron cảm giác xúc giác và đau, xung đi từ cảm biến vào tủy sống, bỏ qua thân tế bào & Cấu tạo (sợi trục chia đôi gần thân) đã được xác lập bằng giải phẫu. Việc dẫn truyền không cần thân kích thích được đã được chỉ ra trên một mô hình đã kiểm chứng của nơ-ron này: không có thân kích thích được, xung vẫn đi qua chỗ rẽ nhánh tin cậy như cũ & \cite{AmirDevor2003} & lý thuyết \\
5 & Một đuôi gai: nơ-ron khứu giác mang một loại thụ thể và truyền một đường đầu vào & Tổng quan các thí nghiệm với gen thụ thể: mỗi nơ-ron khứu giác biểu hiện một gen thụ thể từ một họ lớn & \cite{Mombaerts2004} & đã đo \\
6 & Hai đuôi gai: ở bộ phát hiện hướng âm thanh, đầu vào từ tai trái và tai phải bám vào các đuôi gai khác nhau & Giải phẫu và các bản ghi ở động vật có vú, được tổng hợp trong một tổng quan: đầu vào từ hai tai được tách ra trên hai đuôi gai đối diện & \cite{Grothe2010} & đã đo \\
7 & Tách đầu vào ra hai đuôi gai làm cổng AND chặt hơn & Được chỉ ra trên mô hình: sự bão hòa~\eqref{eq:shunt} trên một đuôi gai không cho các đầu vào từ một phía tới được ngưỡng, và khả năng phát hiện trùng khớp được cải thiện. Chưa có thí nghiệm trực tiếp nào thay đổi cách bố trí đầu vào & \cite{AgmonSnir1998} & lý thuyết \\
8 & Khoảng $7\cdot10^{10}$ nơ-ron \nt{GLUT} nhỏ của mạch học vận động, mỗi nơ-ron $\sim4$ đầu vào & Đếm tế bào bằng phân đoạn đẳng hướng: phần não này ở người có khoảng $69$ tỉ nơ-ron. Ghi một tế bào như vậy ở chuột cống sống: khi chạm, các chùm dòng rời rạc từ vài đầu vào đến, và tế bào phát xung khi chúng cộng lại & \cite{Azevedo2009,Chadderton2004} & đã đo \\
9 & Phần tử ngưỡng có $n$ đầu vào phân tách không quá $P_{\max}\approx2n$ mẫu ngẫu nhiên~\eqref{eq:cover}; với nơ-ron đầu ra của mạch học vận động, giới hạn là $\sim2\cdot10^5$, còn ước tính thực tế là $\sim4\cdot10^4$ phép ánh xạ & Giới hạn là kết quả kinh điển của hình học tổ hợp. Ước tính thực tế là lý thuyết dung lượng với trọng số chỉ dương và có dự trữ cho nhiễu, áp dụng cho số đầu vào đã đo của nơ-ron này & \cite{Cover1965,Brunel2004} & lý thuyết \\
10 & Bộ trộn với ít đầu vào cần để mở rộng đầu vào; con số ``bốn'' gần tối ưu & Lý thuyết: số chiều của biểu diễn do lớp bộ trộn tạo ra đạt cực đại xấp xỉ tại số đầu vào quan sát được trên giải phẫu; giả thuyết mở rộng ban đầu có trong các công trình kinh điển & \cite{Marr1969,Albus1971,LitwinKumar2017} & giả thuyết \\
11 & Cây lớn: $10^4$--$10^5$ đầu vào & Đếm khớp thần kinh trên ảnh hiển vi điện tử: $\approx1.7\cdot10^5$ khớp trên một nơ-ron \nt{GABA} đầu ra của mạch học vận động ở chuột cống; khoảng $30\,000$ đầu vào kích thích và $1\,700$ đầu vào ức chế trên một nơ-ron \nt{GLUT} của hồi hải mã & \cite{NapperHarvey1988,Megias2001} & đã đo \\
12 & Các nhánh của cây lớn là những mạch con riêng với ngưỡng riêng; nơ-ron là một mạng nhiều lớp nhỏ & Kích thích điểm tại hai vị trí trên đuôi gai mảnh: các đầu vào trên cùng một nhánh cộng theo đường sigmoid, trên các nhánh khác nhau thì cộng tuyến tính. Trong đuôi gai của nơ-ron vỏ não người đã tìm thấy các xung canxi phân cấp cho phép một tế bào giải một bài toán không tách được tuyến tính. Mô hình: để tái tạo một nơ-ron cần một mạng sâu & \cite{Polsky2004,Gidon2020,Beniaguev2021} & đã đo \\
13 & Đuôi gai kiêm đầu ra: mỗi nhánh của nơ-ron \nt{GABA} võng mạc là một bộ phát hiện hướng riêng & Ghi canxi hai photon trong từng nhánh: tín hiệu trong nhánh chọn lọc với chuyển động từ thân tế bào ra đầu mút, còn điện áp của thân thì không & \cite{Euler2002} & đã đo \\
14 & Số đầu vào đi theo nhiệm vụ (mệnh đề~\ref{prop:dendrites}) & Cấu tạo các lớp đã được đo (dòng 3--13); việc số đầu vào được chọn vì tốc độ hay vì phân lớp chỉ mới được chỉ ra bằng lý thuyết và mô hình cho một số lớp & \cite{LitwinKumar2017,AgmonSnir1998} & giả thuyết \\
\bottomrule
\end{longtable}}
